\section{Scope, Method, and Qualifications}

\subsection{Reader, question, and governing profiles}

This is Part II of \emph{The Creature Before God}, a discursive series for a serious general reader. Part I, \emph{Ontological Vertigo}, treats the apprehension of asymmetric dependence and the fork of recoil or adoration. This article asks how the response becomes a living sacrifice, how obligation can be freely rendered, how worship takes common form, and how justice passes through the neighbor. The current universal Latin Sunday obligation is a bounded illustration, not the article's principal subject. Part III treats charity and the filial perfection of the return.

The article deliberately combines theology and canon law under the discursive-articles profile. Metaphysics, Scripture, patristic witness, Thomistic moral theology, and authoritative Catholic teaching govern the account of gift, justice, religion, command, grace, hope, and charity. Canon law governs only the identified juridical conclusions. A theological reason does not by itself enact a canon; the existence of a canon does not settle every theological interpretation or personal culpability judgment.

``Due return,'' ``responsive justice,'' the four-term grammar of dependence, duty, precept, and culpability, and the language of ontological vertigo's ongoing resolution are project synthesis. \emph{Religion} is used principally in Aquinas's technical sense: the moral virtue annexed to justice that renders worship to God. It is not used as a synonym for the theological virtues or for every institutional feature of a religious body.

\begin{threestable}{Source class}{Function in the article}{Governing boundary}
Sacred Scripture & Supplies creation and mercy before response, covenant before command, the Decalogue, worship and justice, the unrepayable gift, and the living sacrifice. & Exegetical synthesis remains bounded by the cited passages; no single verse is made a complete theory of law.\\
Patristic witness & Shows divine non-need, command for creaturely good, common Sunday worship, grace-enabled obedience, and self-offering in Christ. & Representative witnesses are not a mechanical consensus count or current legislation.\\
Thomistic theology & Provides justice, religion, unequal due, moral necessity without coercion, devotion, kinds of law, natural and positive determination, grace, hope, and charity's formal role. & School theology and adopted synthesis remain distinct from dogma and current canon law.\\
Catechetical and conciliar teaching & Governs the doctrinal account of moral law, Church precepts, the Decalogue, the Lord's Day, culpability, hope, worship, and common liturgical life. & The \emph{Catechism} is authoritative teaching, not the promulgated text of the Latin Code.\\
Current Latin canon law & Establishes subjects, feast-day duties, satisfaction, impossibility, dispensation, commutation, and territorial competence. & The official Latin text, amendments, jurisdiction, particular law, and facts control; no individual case is adjudicated.\\
Project synthesis & Orders gift, due, precept, fault, hope, and self-offering into Part II of the series. & The terminology and literary movement are not attributed wholesale to any source.\\
\end{threestable}

\Needspace{8\baselineskip}
\subsection{Legal scope and currentness}

\begin{threestable}{Element}{Controlling scope}{Limit}
Governing body of law & 1983 Code of Canon Law as amended; universal law of the Latin Church. & CIC c.~1 excludes applying the Latin Code as such to Eastern Catholics; CCEO c.~881 and proper Eastern law are not analyzed.\\
Authoritative text and translation & Promulgated Latin and applicable later Latin legislation control; Holy See English pages are working reference renderings. & English wording does not displace the official text, canonical tradition, or a later promulgated norm.\\
Persons and territory & CIC cc.~1, 11--13 and 1246--1248; no named territory beyond universal law. & Other holy days, dispensations, and observance can depend upon episcopal-conference, diocesan, proper, or other applicable law.\\
Material facts & Abstract adult Latin Catholic with sufficient use of reason; no serious reason, impossibility, dispensation, commutation, privilege, or special law unless expressly discussed. & Real cases require complete facts and competent pastoral or canonical judgment.\\
Currentness & Official Code pages, the Holy See code and amendment archive, and the Holy See's published compilation of authentic interpretations checked through 20 July 2026. & No located amendment or published authentic interpretation modified CIC cc.~1246--1248 in a way material to this article. This bounded audit is not proof against unpublished or unexamined particular or proper law.\\
Advice boundary & General theological and canonical study. & Not legal or canonical advice, a decree, a culpability judgment, or a substitute for a pastor, confessor, competent authority, or qualified canonist.\\
\end{threestable}

\subsection{Included and excluded scope}

Included are Rom 11:33--12:5, Exod 19--20, Psalm 116, the Decalogue's twofold order to God and neighbor, Irenaeus on divine non-need, Augustine on true sacrifice, Gregory Nazianzen and Chrysostom on gift and mercy, Justin on the unified Sunday assembly, Aquinas on justice, religion, sacrifice, devotion, grace, and hope, Vatican II on Eucharistic and ordinary self-offering, the \emph{Catechism}'s treatment of Church precepts and Sunday, and current universal Latin law in CIC cc.~1, 11, 85--93, 1244--1248, and 912--923 where separately noted.

Excluded are a comprehensive treatise on moral theology, virtue, conscience, sin, sacramental theology, liturgy, ecclesiology, or canonical interpretation; a calendar of territorial holy days; national conference and diocesan norms; Eastern Catholic law beyond naming the boundary; obligations of clerics or religious arising from office, vows, constitutions, or proper law; remote-participation questions; Eucharistic reception in irregular or disputed circumstances; canonical penalties; and any judgment about a named person's obligation, excuse, sin, or imputability.

The Sunday case is illustrative, not exhaustive. The Ten Commandments and the Latin Sunday precept do not stand at the same level of authority. The article does not flatten natural moral law, revealed divine law, and ecclesiastical positive law into one inventory.

\subsection{Material qualifications}

\begin{enumerate}[leftmargin=1.45em,itemsep=0.25em]
\item Creation is neither compelled by a lack in God nor antecedently owed to a recipient not yet existing. Genuine created dues arise within the order willed by divine wisdom (ST I, q.~21, a.~1, ad 3).
\item ``Debt'' and ``return'' are analogical, not commercial. No creature confers actuality, beatitude, or utility upon God or makes God a debtor by worship.
\item Religion is a moral virtue annexed to justice because worship is due but cannot be rendered in absolute equality. It is not itself faith, hope, or charity (ST II--II, q.~81, a.~5).
\item The necessity of command differs from physical or psychological constraint. Obligation does not eliminate voluntariness, while external compliance alone does not constitute complete virtue.
\item Natural reason grounds worship generically; determinate rites and times can arise from divine or human positive law (ST II--II, q.~81, a.~2, ad 3; q.~85, a.~1, ad 1).
\item The Decalogue is revealed and expresses natural moral law. Its moral substance, its covenantal promulgation, and its ceremonial or juridical determinations must be distinguished.
\item Sunday fulfills the Sabbath's spiritual truth in the Paschal mystery; it is not identified with every Mosaic Sabbath determination (CCC 2174--2176).
\item CIC c.~1247 contains both Mass participation and festal-rest duties. The latter is governed by what impedes worship, joy, works of mercy, and appropriate rest, not by a universal ban on all work.
\item Participation in Mass and reception of Holy Communion are distinct juridical and sacramental questions. CIC cc.~912--923 govern reception; this article decides no sacramental-access case.
\item Excuse, impossibility, dispensation, and commutation have different grounds and effects. Private preference is not competent dispensation.
\item Objective grave matter and mortal culpability are distinct. Full knowledge and deliberate consent must be established; a general article cannot judge a person's soul (CCC 1857--1861).
\item Violation of a moral or disciplinary norm does not without an applicable penal norm establish a canonical delict or penalty.
\item Hope is not a species of justice. It trusts divine assistance for the arduous supernatural good; presumption and despair each falsify that trust.
\item Grace does not abolish command or secondary agency. The New Law's inward gift heals and elevates the person to act freely (ST I--II, q.~106, aa.~1--2).
\item Charity perfects and commands the acts of other virtues by ordering them to the final end. It does not make injustice or impiety optional (ST II--II, q.~23, a.~8).
\item The living sacrifice is neither autonomous nor commensurate with Christ's. The Church's self-offering depends wholly upon union with her High Priest and Victim and his perfect prior oblation.
\item Worship is not proposed as therapeutic self-management. Its objective truth and formal fruits do not promise sensible calm, and continuing disturbance is not treated as evidence of refusal.
\item The resolution at issue does not reduce dependence. Justice, religion, and hope bring action into agreement with it; charity's direct unification of desire remains the subject of Part III.
\end{enumerate}

\subsection{Translations, rights, and review}

Scripture is cited by book and locus, with short working-English phrases checked against the Holy See-hosted NAB witness. Patristic working English comes from identified public-domain ANF or NPNF witnesses. Aquinas is cited by work, part, question, article, and reply; Latin loci are governed by the Corpus Thomisticum, with public English used as an aid. Official Catholic and canonical texts retain their own status. No bulk source text, modern commentary, or source-site wrapper is incorporated.

The legal currentness audit and core online witnesses were checked through 20 July 2026; the newly added patristic, conciliar, and Thomistic loci were checked on 21 July 2026. This is a source-audited working study. It has not received independent scriptural, patristic, Thomistic, moral-theological, liturgical, canonical, literary, or ecclesiastical review. It claims no imprimatur, nihil obstat, legal opinion, or ecclesiastical approval.
